% Kweave — formatted for SAGE Publications (sagej class)
\documentclass[Afour,sageh,times]{sagej}

\usepackage[utf8]{inputenc}
\usepackage{moreverb,url}
\usepackage{booktabs}
\usepackage{amsmath}
\usepackage{enumitem}
\usepackage{xcolor}
\usepackage[colorlinks,bookmarksopen,bookmarksnumbered,citecolor=red,urlcolor=red]{hyperref}

\def\volumeyear{2026}

% "Suggested citations" note used at the end of each section (author scaffolding;
% remove before final submission).
\newcommand{\readingnote}[1]{%
  \par\smallskip\noindent\colorbox{gray!12}{%
  \parbox{\dimexpr\linewidth-2\fboxsep}{\footnotesize\textbf{Suggested citations.} #1}}\par}

\begin{document}

\runninghead{Thompson}

\title{\textsc{Kweave}: knowledge workflow --- embeddings, automation, and visual exploration}

\author{Brent Thompson\affilnum{1}}

\affiliation{\affilnum{1}Thompson Solutions Company, USA}

\corrauth{Brent Thompson, Thompson Solutions Company.}

\email{thompsonsolutionscompany@gmail.com}

\begin{abstract}
Decades of curated scientific literature encode domain knowledge that is expensive to
formalize: building an ontology and defining its relationships remains a slow, expert-driven
process that goes well beyond named-entity recognition. We present \textsc{Kweave}
(Knowledge Workflow: Embeddings, Automation, and Visual Exploration), an automated, multi-agent
workflow that turns a curated literature corpus into a standards-compliant knowledge graph.
An orchestrating large language model (LLM) coordinates specialist agents and reproducible
tools that, in one ordered pipeline, discover a shared concept vocabulary from abstracts,
mine each full-text paper for values and exact relational triples, and emit an OWL~2 ontology
with JSON-LD instances interoperable with the Basic Formal Ontology (BFO), the Common Core
Ontologies (CCO), and MDS-Onto. Relation extraction combines an LLM with the REBEL model so
that causal and relational facts are recorded as stated; a LoRA fine-tuning step closes a
learning loop on each finished ontology. Every run produces two deliverables---a GraphDB-ready
repository and a cemento-compatible draw.io diagram---and is gated by structural validation
(reasoner consistency, ontology-pitfall scanning, and upper-ontology alignment). We describe
the system, its evolution across six pipeline generations into a single modular package, and an
evaluation methodology grounded in reproducibility and explainability.
\end{abstract}

\keywords{ontology learning, knowledge graph construction, large language models,
neurosymbolic systems, relation extraction, materials data science, MDS-Onto, BFO, CCO,
provenance, ontology validation, LoRA fine-tuning}

\maketitle

\section{Introduction}
General-purpose language models capture some of the connections latent in published research,
but the majority of their learned representations are of uneven quality for any specific domain,
and they do not by themselves yield an explicit, editable ontology. Knowledge graphs make such
structure explicit, queryable, and reusable~\cite{hogan2021kg,pan2024unifying}. The bottleneck
that motivates this work is the long and labor-intensive process of \emph{ontology creation and
relationship definition}---the step beyond entity recognition and subject coverage that today
requires trained ontologists and coordinated community effort~\cite{smith2007obo}.

The setting that grounds this work is materials data science, where a research group accumulates,
over years, a curated reference library of papers and would like that library to become a queryable,
standards-aligned knowledge base rather than a static reading list. Our corpus is a Zotero group
library organized into per-topic collections; runs to date span gallium-arsenide and
cadmium-telluride (CdTe/CdSeTe) photovoltaics, perovskite solar cells, transmission-electron-microscopy
studies of semiconductors, copper metallization, and electron-microscopy analysis of contact
corrosion---the photovoltaics and materials-durability topics typical of a solar-energy materials
group. Each collection is heterogeneous in vocabulary: the same physical quantity (open-circuit
voltage, contact resistivity, sintering temperature) appears under many surface forms across papers,
which is precisely the normalization problem an ontology is meant to solve.

We address this with \textsc{Kweave}, an \emph{agentic} workflow: LLMs are given tools and run in
a coordinated loop to produce a reproducible and explainable knowledge graph. The system runs in two
modes from the same code path---\emph{unsupervised}, in which the concept vocabulary is discovered
automatically from the corpus's abstracts, and \emph{supervised}, in which a curated concept list is
supplied and the discovery stage is skipped---so a group can bootstrap a new domain or reuse a
hand-tuned schema with a single flag. Our contributions are:
\begin{enumerate}[itemsep=2pt,topsep=2pt]
  \item A single ordered, multi-provider pipeline that converts a curated Zotero corpus into an
        OWL~2 ontology and JSON-LD instances, supporting both unsupervised (auto-discovered
        concepts) and supervised (provided concept list) modes.
  \item A dual extraction stage in which an LLM mines per-concept values while REBEL extracts
        exact subject--predicate--object triples in parallel, capturing relational facts as stated.
  \item Standards grounding to BFO/CCO/MDS-Onto with a mandatory validation gate
        (reasoner consistency, OOPS!\ pitfalls, and alignment coverage).
  \item A LoRA learning loop that fine-tunes on each finished ontology's terms so later runs
        improve on new material.
  \item Two ready-to-use deliverables per run: a GraphDB import repository and a cemento draw.io
        diagram.
\end{enumerate}
\readingnote{Cite a general knowledge-graph reference (Hogan et al., \emph{Knowledge Graphs},
ACM Computing Surveys 2021) and an LLM$+$KG roadmap (Pan et al., IEEE TKDE 2024) to frame the
problem; cite the OBO Foundry (Smith et al., \emph{Nature Biotechnology} 2007) for the cost of
manual, coordinated ontology engineering.}

\section{Background and Related Work}
\textbf{Foundational ontologies.} The Basic Formal Ontology (BFO) is a top-level ontology adopted
as ISO/IEC~21838-2~\cite{bfo}; the Common Core Ontologies (CCO) extend it at the mid level with
reusable classes for information, artifacts, and processes; and MDS-Onto---the Materials Data Science
ontology maintained by the Case Western Solar Durability and Lifetime Extension (SDLE) center---supplies
the domain vocabulary, published under the canonical namespace
\texttt{https://cwrusdle.bitbucket.io/mds/}. Materials-specific upper resources such as the
Platform MaterialDigital (PMD) core ontology and the QUDT vocabulary for quantities and units occupy
the same grounding layer. Anchoring generated classes in these shared resources is what makes the
output interoperable and reusable rather than a private schema, following the design discipline
established by the OBO Foundry~\cite{smith2007obo}; it is also the ``semantic grounding'' requirement
that distinguishes a trustworthy materials knowledge graph from a free-text extraction dump.

\textbf{Automated knowledge extraction in materials science.} Turning unstructured scientific text
into structured, machine-actionable data---synthesis recipes, process parameters, device metrics, and
the relationships among them---is a long-standing goal that NLP and, more recently, LLMs have made
tractable. Classical ontology-learning systems such as Text2Onto established the task before the neural
era; transformer keyphrase methods (e.g.\ KeyBERT) and, now, instruction-tuned LLMs have raised both
recall and the fidelity of the extracted structure. \textsc{Kweave} sits in this line but targets the
harder end of it: not only \emph{which} concepts a corpus discusses, but their values, their exact
stated relations, and an editable OWL~2 ontology that ties them to an upper-level standard.

\textbf{LLM-based knowledge-graph construction.} Recent surveys map a fast-moving space of
LLM-driven extraction and fusion methods~\cite{kgsurvey,llmonto,pan2024unifying}. NeOn-GPT is the
closest prior pipeline, translating natural-language descriptions into Turtle ontologies, and its
reported limits---LLMs lacking the procedural and reasoning rigor for full ontology
construction---directly motivate our tool-augmented, validated design~\cite{neongpt}. Related
extraction-from-text systems include schema-driven prompting approaches such as
SPIRES/OntoGPT~\cite{caufield2024ontogpt}, which likewise constrain an LLM to a target schema. Our
response to the rigor gap is architectural: the LLM never emits the ontology unchecked, because a
structural validation gate (reasoner consistency, pitfall scanning, and upper-ontology alignment)
stands between generation and publication.

\textbf{Neurosymbolic agents, relation extraction, and verification.} \textsc{Kweave} is a
neurosymbolic system in the literal sense: a neural extraction layer (the LLM plus REBEL) feeds a
symbolic layer (an OWL~2 ontology with a description-logic reasoner and shape checks). REBEL casts
relation extraction as sequence generation, emitting triplets end-to-end~\cite{rebel}; we use it as a
reproducible, separately-versioned tool alongside the LLM so that relational and causal facts are
captured as stated in the text rather than paraphrased. The orchestration follows the conversational,
tool-using paradigm of AutoGen~\cite{autogen}, building on reasoning-and-acting and self-tooling
methods such as ReAct~\cite{yao2023react} and Toolformer~\cite{schick2023toolformer}, with the LLM as
planner and the neural models as callable components. Where many agentic systems treat the model's
output as authoritative, we attach provenance to every assertion and gate publication on verification,
which is the guardrail an autonomous literature-to-KG assistant needs to be trusted in a research
workflow.
\readingnote{Ontologies: add CCO and an MDS-Onto/SDLE reference for your group's vocabulary.
Ontology learning: Text2Onto (Cimiano \& V\"olker, 2005) and OntoGPT/SPIRES (Caufield et al.,
\emph{Bioinformatics} 2024). Agents/tool use: ReAct (Yao et al., 2023), Toolformer (Schick et al.,
2023), and a recent multi-agent survey to complement AutoGen.}

\section{System Architecture}
The workflow is a single linear pipeline with two hard ordering invariants: JSON-LD instances and
the LoRA fine-tune run \emph{only after} the ontology is built. The LLM acts as orchestrator and
reasoner; the neural models (REBEL, the fine-tuned extractor) are versioned, seeded, logged tools.
Table~\ref{tab:stages} summarizes the stages.

\begin{table*}[h]
\small\sf\centering
\caption{The ordered \textsc{Kweave} pipeline. Steps 5, 5b, and 6 execute only after Step~4.\label{tab:stages}}
\begin{tabular}{@{}llp{8.6cm}@{}}
\toprule
\textbf{Step} & \textbf{Name} & \textbf{Function} \\
\midrule
0 & Input        & Fetch papers + PDF text from Zotero; choose supervised/unsupervised mode \\
1 & Concepts     & Extract concepts from abstracts; normalize to a shared 30--80 term list \\
2 & Mine         & Per paper: LLM extracts value+quote per concept; REBEL extracts triples \\
3 & Consolidate  & Reconcile concepts and triples into the final ontology terms \\
4 & Ontology     & Build OWL~2 TTL grounded in MDS/CCO/BFO; \emph{validation gate} \\
5 & JSON-LD      & Emit per-paper and combined JSON-LD instances (GraphDB repo) \\
5b& Diagram      & Tag concepts (study stage / supply chain); build cemento draw.io \\
6 & LoRA         & Fine-tune adapter on the final ontology terms (terminal step) \\
\bottomrule
\end{tabular}
\end{table*}

\noindent
\textbf{Provider-agnostic agents.} Each agent is configured against any OpenAI-compatible
endpoint, so the same code runs on a hosted frontier model (the default is \texttt{claude-sonnet-4-6})
or a local 7--8B-parameter Llama~\cite{touvron2023llama} served through an OpenAI-compatible shim,
with only an environment variable changed. Structured outputs are obtained through typed result
schemas rather than free-text parsing: three extraction agents declare Pydantic return types---an
extractor that emits per-abstract concepts as \texttt{(canonical label, paper term, relevance)}
triples, a normalizer that collapses the corpus's raw labels into a deduplicated set of 30--80
ontology-ready terms, and a miner that returns, for each concept, the paper's exact value and the
single supporting sentence. The tool/function schema is generated from these types, so the model is
constrained to well-formed output. For grounding, a \emph{fixed-choice} taxonomy is enforced at the
type level (the study-stage and supply-chain tags are constrained \texttt{Literal} enumerations),
which is what makes the smaller local models usable for tagging.

\textbf{Tunable parameters.} A handful of pinned knobs make runs reproducible and tractable:
concept extraction keeps the top 25 concepts per abstract; full-text mining truncates each paper to
80{,}000 characters; tagging batches 40 concepts per request; and a 0.5\,s inter-call delay respects
provider rate limits. Tool choice can be forced on so that every structured call returns through the
typed schema. Because processing is serial and seeded, two runs over the same corpus with the same
model produce the same artifacts.

\textbf{Grounding and class assignment.} The emitter places every domain class beneath one of eight
MDS-Onto branches---Measurement, Material, Process, Device, Characterization, Reliability, Sample, and
Economics---all subclassing a common \texttt{mds:Concept}, with Reliability modeled as a subclass of
Characterization. A keyword-rule classifier assigns each discovered concept to its branch, and each
concept becomes an OWL \texttt{DatatypeProperty} on a domain \texttt{SchemaRecord} class (itself a
subclass of \texttt{mds:ResearchPublication}), carrying an \texttt{rdfs:label}, a \texttt{skos:broader}
link to its branch, and up to five \texttt{skos:example} values drawn from the actual paper terms. A
central namespace registry fixes one canonical MDS IRI (the slash form) and routes each domain to its
own sub-namespace---for example solar cells, GaAs, electron microscopy, and copper metallization each
resolve under their MDS sub-path---so IRIs join cleanly across runs.

\textbf{Relation extraction as a versioned tool.} REBEL (\texttt{Babelscape/rebel-large}) is invoked
during mining with beam search and the model's canonical special-token decoding to recover flat
subject--predicate--object triples; its model name and revision are pinned for reproducibility, and the
stage degrades to a safe no-op if the dependency or weights are unavailable, so the rest of the pipeline
still completes. Extracted triples are persisted into the GraphDB repository as both CSV and JSON-LD.

\textbf{Explainability and reproducibility.} Model versions and seeds are pinned; each assertion
carries a provenance record---source paper, text span, tool, model and version, seed, and
confidence---modeled on the W3C PROV vocabulary~\cite{provo}; REBEL triples are exported with a
\texttt{prov:wasDerivedFrom} link to their source paper, and per-paper JSON-LD instances retain an
\texttt{owl:sameAs} link to the article DOI. The validation gate must pass before the repository is
published.
\readingnote{Structured output: cite the tool/function-calling mechanism you rely on (e.g.\
constrained/grammar-based decoding) and the PydanticAI/Instructor libraries. Provenance: PROV-O
(W3C Recommendation, 2013). Local models: LLaMA (Touvron et al., 2023).}

\section{Implementation and Evolution}
\textsc{Kweave} is implemented as a single modular Python package (\texttt{kw}) exposing one ordered
runner, with one module per responsibility---\texttt{zotero} (input), \texttt{extract} and
\texttt{rebel} (mining), \texttt{taxonomy} and \texttt{tagger} (grounding vocabulary), \texttt{ontology}
(OWL~2 / JSON-LD emission), \texttt{validate} (the gate), \texttt{drawio} (diagram), \texttt{lora}
(learning loop), with \texttt{config}, \texttt{models}, \texttt{store}, and \texttt{pipeline} holding
the namespace registry, the typed data contract, serialization, and the runner respectively. This flat
layout is the deliberate consolidation of an earlier sprawl of overlapping scripts: the design rule was
\emph{one owner per step, no unnecessary complexity}. Papers and full text are retrieved through the
Zotero API (group library, per-collection); the ontology and JSON-LD are produced by a Turtle/JSON-LD
emitter that classifies concepts into the eight MDS branches above; output conforms to
OWL~2~\cite{owl2} and JSON-LD, and structural constraints can be expressed in SHACL~\cite{shacl}.
Between stages, a wide-format schema CSV (one row per paper, one column per concept, each cell packing
the extracted value with its supporting quote) serves as the on-disk handoff the emitter consumes.

\textbf{Grounding tags and the diagram.} As part of grounding, each concept is tagged with an MDS-Onto
\emph{study stage}---one of twelve values spanning the experimental lifecycle (synthesis, formulation,
materials processing, sample, tool, recipe, data, data processing, result, analysis, modeling, and
results-and-metadata)---and one or more of six \emph{supply-chain levels} (materials, subcomponent,
component, assembly, subsystem, system). These taxonomies live in a single source module so the tagger
and the diagram can never diverge. The optional draw.io deliverable lays the tagged concepts out as
colour-coded swimlanes by study stage, with a central workspace for an ontologist to drag and connect
nodes, and embeds the MDS-Onto and Cemento shape libraries as palette pages so the generated map is
immediately editable in the cemento ontology-drawing tool.

\textbf{The learning loop.} After a domain ontology is finalized, a Low-Rank Adaptation
(LoRA)~\cite{hu2021lora} adapter is fine-tuned on the run's final ontology terms and mined
text$\rightarrow$value pairs; quantized variants such as QLoRA~\cite{dettmers2023qlora} make this
feasible on modest hardware. The supervised dataset pairs each paper's text with the JSON object of its
extracted concept values, plus a vocabulary example that teaches the canonical term set itself; training
uses a small-rank adapter (rank~16, applied to the attention query and value projections) over a few
epochs. The step degrades gracefully---absent a GPU or the training dependencies it still writes the
dataset and a manifest, so the pipeline always completes and training can be resumed later. The intent
is that successive runs reuse prior ontologies to ingest new material more accurately.

The current design is the consolidation of six pipeline generations
(Table~\ref{tab:history}). Early versions explored classical NLP and transformer keyword
extraction; mid versions introduced structured LLM extraction and two-stage discovery with
full-text mining; later versions added provider independence and concept enrichment. The work
described here unifies these---together with a previously separate ontology/JSON-LD emitter and a
draw.io generator---into one deduplicated pipeline with a single data contract and a validation
gate.

\begin{table}[h]
\small\sf\centering
\caption{Evolution of the pipeline into \textsc{Kweave}.\label{tab:history}}
\begin{tabular}{@{}lll@{}}
\toprule
\textbf{Gen.} & \textbf{Method} & \textbf{Key advance} \\
\midrule
V1 & spaCy + TF-IDF        & offline concept counts (baseline) \\
V2 & KeyBERT               & semantic keyphrases \\
V3 & LLM structured        & per-paper typed extraction \\
V4 & Two-stage + full text & concept discovery + source quotes \\
V5 & Provider-agnostic     & any OpenAI-compatible model \\
V6 & Enrichment            & MDS tagging + concept-map diagrams \\
\midrule
\multicolumn{3}{@{}l}{\emph{This work} (\textsc{Kweave}): unified modular package;} \\
\multicolumn{3}{@{}l}{OWL~2/JSON-LD + diagram; REBEL + LoRA; validation.} \\
\bottomrule
\end{tabular}
\end{table}
\readingnote{Standards: OWL~2 (W3C, 2012), JSON-LD (W3C, 2020), SHACL (W3C, 2017). Tooling:
ROBOT (Jackson et al., 2019) for reasoning/reporting; the GraphDB and Zotero documentation as
software references. Fine-tuning: LoRA (Hu et al., 2021), QLoRA (Dettmers et al., 2023), and an
instruction-tuning reference (e.g.\ Wei et al., FLAN, 2022).}

\section{Evaluation}
Because manually constructing a reference ontology for every domain is impractical---and the cost
and difficulty of manual ontology engineering is well established in the literature---we evaluate
the \emph{generated artifact} with structural metrics rather than a human study:
\begin{itemize}[itemsep=2pt,topsep=2pt]
  \item \textbf{Logical consistency}: the OWL~2 ontology passes a DL reasoner (e.g.\ HermiT or
        ELK, invoked via ROBOT) with no unsatisfiable classes. The reasoner and pitfall checks are
        exposed behind a stable interface in the validation module so that the external tools can be
        wired in without changing the pipeline.
  \item \textbf{Modeling quality}: an ontology-pitfall scan (OOPS!~\cite{oops2014}) reports no
        critical pitfalls.
  \item \textbf{Interoperability}: the alignment ratio---the fraction of declared classes that
        subclass a recognized upper or mid-level parent (BFO, CCO, MDS-Onto, or PMD). The gate
        requires this ratio to meet a fixed threshold (0.80) in addition to passing the reasoner and
        pitfall checks; a run that fails is not published.
\end{itemize}
This artifact-centric protocol is itself the verification guardrail the system needs to be trusted as
an autonomous extractor: the ontology is published only if it is logically consistent, free of critical
modeling pitfalls, and demonstrably grounded in a shared standard. On a hand-built reference produced by
an earlier run (a gallium-arsenide photovoltaics ontology), all seven declared classes aligned to an
upper parent (alignment ratio 1.0), indicating the grounding step behaves as intended. The pipeline has
been exercised across the photovoltaics and materials-durability domains in the corpus---GaAs and CdTe
solar cells, perovskites, TEM semiconductor studies, copper metallization, and electron-microscopy
contact corrosion---and a full multi-domain evaluation, including runs driven by a local 7--8B model, is
in progress. Reproducibility is supported by pinned model and seed values and by the provenance attached
to every assertion.
\readingnote{Ontology evaluation: OOPS! (Poveda-Vill\'on et al., 2014) and a general ontology
evaluation reference (e.g.\ Vrande\v{c}i\'c, 2009; or OntoMetrics). Reasoners: HermiT (Glimm et al.,
2014) and ELK (Kazakov et al., 2014). If you add a human-quality study later, cite an
inter-annotator agreement methodology.}

\section{Discussion and Limitations}
The architecture deliberately favors simplicity: a linear pipeline with one owner per step, no
speculative agent loop. This is a considered position---an earlier incarnation of the project had
fragmented into several entry points with duplicated tagging and emission logic and no shared data
contract, and the consolidation into one ordered runner with a single typed contract is what made an
end-to-end, validated run possible at all. The same discipline shapes how the neural tools are used:
REBEL and the reasoner are optional, separately-versioned components that degrade to safe no-ops when
absent, so a run never fails because an extra model is missing, at the cost of those signals being
unavailable in that run.

Open challenges remain. (i)~Strict structured output from small local models is not always reliable;
we mitigate it with type-constrained schemas and conservative concept counts, but a local-model run can
still under-fill the schema relative to a frontier model. (ii)~Entity resolution at the merge step is
deliberately shallow: REBEL triples and LLM concepts are persisted side by side, and collapsing the
same real-world entity named differently in each stream into a single node is future work rather than a
solved problem. (iii)~The supervision available to the LoRA loop is modest---it seeds from the run's own
ontology terms and mined pairs plus prior runs---so the compounding-improvement claim is a design
hypothesis we have built the machinery to test, not yet a measured result. (iv)~On engineering
robustness, processing is serial (on the order of several minutes for tens of papers in the mining
stage) and lacks per-paper checkpointing and retry/backoff, so a mid-run provider error currently
restarts the run; this is the most immediate hardening target. Finally, structural validation
establishes internal soundness---consistency, modeling hygiene, and grounding---but not domain
correctness; whether a class genuinely captures the materials concept it names still benefits from
expert review, and the editable draw.io diagram exists precisely to make that review fast.
\readingnote{Entity resolution/linking: a neural entity-linking reference (e.g.\ Wu et al., BLINK,
2020). Reliability of LLM output: a survey on hallucination/factuality in LLMs. Small-model
structured output: constrained-decoding references.}

\section{Conclusion and Future Work}
We described \textsc{Kweave}, an automated, neurosymbolic workflow that converts a curated literature
corpus into a validated, standards-grounded knowledge graph and an editable diagram, with relational
facts captured by REBEL and a per-ontology LoRA learning loop. By pairing a neural extraction layer with
a symbolic ontology, attaching provenance to every assertion, and gating publication on a structural
verification step, the system aims at the kind of trustworthy, literature-to-knowledge-graph assistant
that semantic materials science needs: one that a research group can point at a Zotero collection and
trust to produce an interoperable, inspectable artifact rather than an unverified extraction.

Future work proceeds along three lines. On \emph{trust and verification}, we will wire the reasoner,
OOPS!, and SHACL checks fully into the gate and report alignment, consistency, and pitfall metrics
across every domain in the corpus, complemented by expert review of a sampled subset. On
\emph{extraction quality}, we will deepen entity resolution at the merge step so that REBEL relations
and LLM concepts collapse onto shared URIs, and add BFO/CCO parents wherever MDS-Onto lacks one. On
\emph{the learning loop}, we will use the trained LoRA adapter directly inside the extraction agents so
that each finished ontology measurably improves the next run, turning the per-ontology fine-tune from a
terminal step into a compounding one. Hardening the runner with per-paper checkpointing and
retry/backoff rounds out the path from research prototype to a dependable autonomous assistant.

\begin{acks}
The author thanks the Case Western SDLE research community for the MDS-Onto vocabulary and the
curated photovoltaics and materials-durability literature that motivated this work.
\end{acks}

\begin{dci}
The author declares that there is no conflict of interest.
\end{dci}

\begin{funding}
The author received no financial support for the research, authorship, and/or publication of this
article.
\end{funding}

\begin{thebibliography}{99}

\bibitem[Arp et al.(2015)]{bfo}
Arp R, Smith B and Spear AD (2015) \textit{Building Ontologies with Basic Formal Ontology}.
Cambridge, MA: MIT Press.

\bibitem[Caufield et al.(2024)]{caufield2024ontogpt}
Caufield JH et al. (2024) Structured prompt interrogation and recursive extraction of semantics
(SPIRES). \textit{Bioinformatics} 40(3).

\bibitem[Dettmers et al.(2023)]{dettmers2023qlora}
Dettmers T et al. (2023) QLoRA: Efficient finetuning of quantized LLMs. In: \textit{Advances in
Neural Information Processing Systems (NeurIPS)}.

\bibitem[Fathallah et al.(2024)]{neongpt}
Fathallah N, Das A, De Giorgis S, Poltronieri A, Haase P and Kovriguina L (2024) NeOn-GPT: A large
language model-powered pipeline for ontology learning. In: \textit{ESWC 2024 Satellite Events}.
Springer, pp.\ 36--50.

\bibitem[Hogan et al.(2021)]{hogan2021kg}
Hogan A et al. (2021) Knowledge graphs. \textit{ACM Computing Surveys} 54(4): 1--37.

\bibitem[Hu et al.(2022)]{hu2021lora}
Hu EJ et al. (2022) LoRA: Low-rank adaptation of large language models. In: \textit{International
Conference on Learning Representations (ICLR)}. arXiv:2106.09685.

\bibitem[Huguet Cabot and Navigli(2021)]{rebel}
Huguet Cabot PL and Navigli R (2021) REBEL: Relation extraction by end-to-end language generation.
In: \textit{Findings of the ACL: EMNLP 2021}, pp.\ 2370--2381.

\bibitem[Knublauch and Kontokostas(2017)]{shacl}
Knublauch H and Kontokostas D (2017) \textit{Shapes Constraint Language (SHACL)}. W3C Recommendation.

\bibitem[LLM-KG Survey(2025)]{kgsurvey}
LLM-empowered knowledge graph construction: A survey (2025). arXiv:2510.20345.

\bibitem[LLM Ontology Engineering Review(2025)]{llmonto}
Large language models for ontology engineering: A systematic literature review (2025).
\textit{Semantic Web Journal}.

\bibitem[Pan et al.(2024)]{pan2024unifying}
Pan S et al. (2024) Unifying large language models and knowledge graphs: A roadmap. \textit{IEEE
Transactions on Knowledge and Data Engineering}.

\bibitem[Poveda-Villal\'on et al.(2014)]{oops2014}
Poveda-Villal\'on M, G\'omez-P\'erez A and Su\'arez-Figueroa MC (2014) OOPS! (OntOlogy Pitfall
Scanner!). \textit{International Journal on Semantic Web and Information Systems} 10(2): 7--34.

\bibitem[Schick et al.(2023)]{schick2023toolformer}
Schick T et al. (2023) Toolformer: Language models can teach themselves to use tools. arXiv:2302.04761.

\bibitem[Smith et al.(2007)]{smith2007obo}
Smith B et al. (2007) The OBO Foundry: Coordinated evolution of ontologies to support biomedical
data integration. \textit{Nature Biotechnology} 25: 1251--1255.

\bibitem[Touvron et al.(2023)]{touvron2023llama}
Touvron H et al. (2023) LLaMA: Open and efficient foundation language models. arXiv:2302.13971.

\bibitem[W3C OWL Working Group(2012)]{owl2}
W3C OWL Working Group (2012) \textit{OWL~2 Web Ontology Language Document Overview}, 2nd edn.
W3C Recommendation.

\bibitem[World Wide Web Consortium(2013)]{provo}
World Wide Web Consortium (2013) \textit{PROV-O: The PROV Ontology}. W3C Recommendation.

\bibitem[Wu et al.(2023)]{autogen}
Wu Q et al. (2023) AutoGen: Enabling next-gen LLM applications via multi-agent conversation
framework. arXiv:2308.08155.

\bibitem[Yao et al.(2023)]{yao2023react}
Yao S et al. (2023) ReAct: Synergizing reasoning and acting in language models. In:
\textit{International Conference on Learning Representations (ICLR)}.

\end{thebibliography}

\end{document}
